\documentclass[11pt,compress,t,notes=noshow, xcolor=table]{beamer}

\input{../../style/preamble}
\input{../../latex-math/basic-math}
\input{../../latex-math/basic-ml}
\input{../../latex-math/optim-basics}

\newcommand{\alphav}{\bm{\alpha}}

\title{Optimization in Machine Learning}

\begin{document}

\titlemeta{
Constrained Optimization
}{
Primal Versus Dual Methods
}{
figure/primal_dual_views.pdf
}{
\item Compare primal and dual viewpoints for the same constrained problem
\item Understand why dual methods are not automatically superior or required for kernels
\item Relate computational cost to the number of variables $d$ and data points $n$
\item See why approximate primal solutions are often easier to interpret and use
}

\begin{framei}[fs=small]{Two Viewpoints for the Same Problem}
\item The primal solves the original constrained problem
$$
\min_{\wv \in \S} F(\wv),
$$
directly in the model parameters $\wv$.
\item The dual solves
$$
\max_{\alphav \ge \zero} L(\alphav),
\qquad
L(\alphav)=\min_{\wv} H(\wv,\alphav),
$$
in the multiplier variables $\alphav$.
\item Thus, primal and dual methods come from the same Lagrangian, but they optimize in different coordinates.
\item Primal iterates live in the original parameter space; dual iterates live in multiplier space and usually need a recovery step back to $\wv$.
\item The main practical differences are dimension, constraint structure, computational cost, and interpretability.
\end{framei}


\begin{framei}[fs=small]{Duality Is a Reformulation}
\item The dual is created from the same Lagrangian as before:
$$
H(\wv,\alphav) = F(\wv) + \sum_{i=1}^m \alpha_i g_i(\wv),
\qquad
\alphav \ge \zero.
$$
\item Passing from primal to dual does \textbf{not} remove optimization difficulty; it changes the coordinates of the problem.
\item In practice, deriving the dual often introduces one multiplier per constraint, extra sign or box constraints on these multipliers, and a mapping back from $\alphav$ to $\wv$.
\item Hence the relevant question is not \enquote{primal or dual in general?}
\item The relevant question is:
\oneliner{Which representation best matches the dimension, constraints, and approximation needs of the current problem?}
\end{framei}


\begin{framev}[fs=footnotesize]{Kernel Methods: Dual Helpful, Not Essential}
\splitVTT[0.53]{
\textbf{Linear SVM in the primal}
$$
\min_{\wv}
\frac{\lambda}{2}\|\wv\|_2^2
+
\sum_{i=1}^n \max\left\{1 - y^{(i)} {\xv^{(i)}}^\top \wv, 0\right\}
$$

\textbf{Soft-margin SVM in the dual}
$$
\max_{\alphav}
\quad
\sum_{i=1}^n \alpha_i
-
\frac{1}{2}
\sum_{i=1}^n \sum_{j=1}^n
\alpha_i \alpha_j y^{(i)} y^{(j)} k_{ij}
$$
\vspace{-0.2cm}
$$
\text{s.t. }
\quad
0 \le \alpha_i \le C,
\qquad
\sum_{i=1}^n \alpha_i y^{(i)} = 0
$$
}{
\begin{itemize}
\item The dual depends only on similarities $k_{ij}$, so replacing dot products by a kernel is immediate.
\item But the primal SVM is already a valid unconstrained optimization problem once the hinge loss is written inside the objective.
\item Therefore, kernel methods require similarities, not duality itself; via the representer viewpoint, many regularized primal solutions can also be written in similarity form.
\item Duality is useful here, but it is \textbf{not} the only reasonable way to solve or kernelize the model.
\end{itemize}
}
\end{framev}


\begin{framev}[fs=footnotesize]{Which Side Is Computationally Cheaper?}
\splitVTT[0.49]{
\textbf{Primal:} $\Xmat^\top \Xmat \in \R^{d \times d}$

\textbf{Dual:} $\mathbf{K} = \Xmat \Xmat^\top \in \R^{n \times n}$

\begin{itemize}
\item The primal typically works with a \textbf{scatter matrix} of size $d \times d$.
\item The dual typically works with a \textbf{similarity matrix} of size $n \times n$.
\item Core storage and linear-algebra cost are therefore often $\order(d^2)$ versus $\order(n^2)$.
\item If $d < n$, primal methods are often cheaper.
\item If $n < d$, dual methods can be cheaper.
\item For very large $n$, storing $\mathbf{K}$ can already be the bottleneck.
\end{itemize}
}{
\imageC[0.98]{figure/primal_dual_regimes.pdf}
}
\end{framev}


\begin{framei}[fs=small]{Approximate Solutions Often Favor the Primal}
\item For convex problems with strong duality, the \textbf{optimal values} satisfy
$$
p^\ast = d^\ast.
$$
\item But this equality holds at optimality; it does \textbf{not} say that an approximately optimal dual iterate already gives a good primal model.
\item In the SVM, the primal weight vector is recovered from the dual variables via
$$
\wv(\alphav) = \sum_{i=1}^n \alpha_i y^{(i)} \xv^{(i)}.
$$
\item A dual iterate with a good objective value can still translate to a poor primal model if the primal-dual relations are not yet accurate enough.
\item By contrast, primal methods directly optimize the objective we actually care about in the original variables.
\item This makes early stopping and numerical approximation easier to interpret in the primal.
\end{framei}


\begin{framev}[fs=small]{Interpretability of Intermediate Iterates}
\splitVTT[0.50]{
\textbf{Primal iterates}
\begin{itemize}
\item are coefficient vectors, probabilities, portfolio weights, or separating hyperplanes
\item can be inspected, constrained, or used directly during optimization
\item are often meaningful even before full convergence
\end{itemize}
}{
\textbf{Dual iterates}
\begin{itemize}
\item are multipliers, shadow prices, or point weights
\item are useful diagnostically, but are one step removed from the final model
\item usually need a primal recovery step before interpretation or deployment
\end{itemize}

\vspace{0.2cm}
\oneliner{If computation stops early, a current primal iterate is usually easier to justify and communicate.}
}
\end{framev}


\begin{framei}[fs=small]{When Primal Methods Are the Default Choice}
\item Primal methods are usually attractive when projections, coordinate slices, or barrier terms are simple enough to compute.
\item They are especially natural when $d$ is smaller than or comparable to $n$, so the feature-space representation is compact.
\item They are often preferable when we want interpretable parameters, coefficient bounds, sparsity patterns, or direct control over feasibility.
\item They are also a good default when the primal can be written as a simple unconstrained or lightly constrained objective, as in hinge-loss or regularized risk formulations.
\item In these settings, there is often little reason to pay the extra indirection of a dual reformulation.
\end{framei}


\begin{framei}[fs=small]{When Dual Methods Can Still Be Attractive}
\item Dual methods become appealing when the number of points $n$ is much smaller than the number of variables $d$.
\item They can be convenient when the dual constraints are structurally simpler than the original primal constraints, often box constraints plus a few equalities.
\item They are also useful when the model is naturally expressed through similarities or kernel evaluations and the matrix $\mathbf{K}$ is still manageable.
\item The dual objective provides lower bounds, and the multipliers themselves can carry information about active constraints or influential observations.
\item Thus, dual methods remain valuable, but only when the reformulation clearly exploits useful structure.
\end{framei}


\begin{framev}{Take-Away}
\begin{enumerate}
\item Primal and dual methods are two views of the same constrained optimization problem.
\item Duality is powerful, but it is not automatically better and it is not required for kernel methods.
\item Computationally, the choice is often driven by dimension: $d < n$ tends to favor primal methods, while $n < d$ can favor dual methods.
\item For approximate solutions, primal iterates are usually easier to interpret, evaluate, and stop early.
\end{enumerate}

\vfill
\oneliner{Practical rule: use a primal method whenever it is simple enough; move to the dual only when the reformulation clearly buys something.}
\end{framev}

\endlecture
\end{document}
